\subsection{How much of this duplicates the libraries?}
\label{sec:dup}

It is natural to suspect that a development of this size, written quickly against a large
library, consists largely of re-proved library facts. We tested the suspicion, and it is
false, but the measurement is worth reporting because the way in which it is false is the
interesting part.

\paragraph{Name overlap.} Extracting every \lean{theorem}, \lean{lemma}, \lean{def},
\lean{abbrev} and \lean{instance} name from the three trees gives 43{,}101 declarations
here, 170{,}982 in \Mathlib{} (at the pinned commit) and 237{,}302 in \TauCeti{}. The count
for this repository was taken before the unused-lemma removal of
Section~\ref{sec:unused}; that removal deletes 6{,}328 declaration names, and repeating the
two exact-match rows on the pruned tree gives 422 (1.1\%) and 513 (1.4\%)---the shares move
by a tenth of a percentage point, which is the size of the effect that a 12\% reduction in
code has on this measurement.

\begin{center}\small
\begin{tabular}{@{}lrr@{}}
\toprule
comparison & declarations & share of the development \\
\midrule
exact name also in \Mathlib{} & 438 & 1.0\% \\
exact name also in \TauCeti{} & 531 & 1.2\% \\
base name also in \Mathlib{} & 829 & 1.9\% \\
base name also in \TauCeti{} & 899 & 2.1\% \\
base name in either & 1{,}159 & 2.7\% \\
\bottomrule
\end{tabular}
\end{center}

Name matching is a weak proxy and it cuts both ways: the project names things
descriptively (\lean{jetPolynomial}, \lean{majorant\_mul}, \lean{rising\_mono}), so semantic
duplication is under-counted.

\paragraph{Where the duplication actually is.} The infrastructure layer is the answer. Of
the 2{,}588 files, eighteen---3{,}855 lines, 0.7\% of the lines---rebuild
standard analysis: smooth cutoffs with all-order bounds, the exponential edge
$\exp(-c/x^2)$, zero extension across it, finite-order jet bookkeeping, the Banach space
$C^k([-1,1];E)$, differentiation under the integral sign (twice, in incompatible forms),
Schwartz density, $L^2$ mollification, word-lattice Sobolev spaces, and the derivative gain
of the heat kernel. That is where \TauCeti{} overlaps, and that is the entire target set for
a library migration. The remaining 99.4\% of the code has no counterpart in any library,
because it is the proof: the profiles, the dyadic charts, the reserved patches, the pulse
covariances, the correction stages, the packet frames.

\paragraph{Duplication of the development with itself.} A different and more serious
question is internal duplication, and here the answer is less comfortable: 74 modules are
unreachable from any entry point (Section~\ref{sec:cleanup}), eleven of them literal copies
of other modules placed in a separate namespace so that a heavy proof could be re-elaborated
under different heartbeat limits, and the same argument is proved twice in
\file{SmoothParameterIntegral} and \file{CompactParameterIntegral} because the two files
were written independently. Internal duplication, not library duplication, is the real
technical debt here.

\paragraph{Measuring it properly.} The right instrument is a type-level comparison, not a
name-level one: \Mathlib{} ships \file{Mathlib/Tactic/DuplicateDecls.lean} (Gerbscheid),
which normalises every theorem type by erasing universes and sorting binders, then hashes
the result, so that two theorems count as duplicates when their statements agree up to
binder order, binder information and universe levels. It cannot be used as shipped---its
declarations are not \lean{public} in the module system, so an importing file cannot call
them---and we vendored the forty lines of its core into a probe.

We validated the probe before trusting it, and the validation is the reason the numbers
below are worth reporting. Two controls were placed in a separate module: \lean{plantedA}
and \lean{plantedB}, the theorems $n=n$ and $m=m$, which duplicate each other; and
\lean{plantedAddComm}, stating $a+b=b+a$, which duplicates \Mathlib's \lean{Nat.add\_comm}.
The probe found exactly two groups involving a project declaration: \lean{[plantedA,
plantedB]} and \lean{[Nat.add\_comm, plantedAddComm]}. It sees what it claims to see.

Running it over 22 of the repository's modules---the infrastructure layer and the first
\file{NavierStokes} files, i.e.\ deliberately sampling where duplication with a library is
most likely---gives:

\begin{center}\small
\begin{tabular}{@{}lr@{}}
\toprule
theorem-duplicate groups in the environment & 1{,}533 \\
project theorems passing the filters & 802 \\
groups involving a project theorem & 2 (both planted controls) \\
groups involving a project theorem, controls removed & \textbf{0} \\
\bottomrule
\end{tabular}
\end{center}

The two counts above the line are lower than in the first measurement of this tree (2{,}401
groups, 944 theorems) because the unused-lemma removal of Section~\ref{sec:unused} took
142 theorems out of the sample as well (6{,}330 out of the tree); the number that matters,
the last row, is unchanged at zero.

So not a single theorem in the sample restates a library theorem, even up to binder
reordering. This does not contradict the infrastructure finding above: the duplication that
exists is duplication of \emph{concepts}, not of statements. \lean{SmoothCutoffs.one\_le\_growth}
is not any \Mathlib{} lemma; it is a bespoke weighted bound whose role a library would
express differently. The cost of that is real, but it is not measured by counting shared
theorems, and the migration effort should be directed by the concept-level audit of
Section~\ref{sec:replace}, not by a duplicate-declaration report.

\paragraph{The probe as a regression check.} Both the probe and its controls are kept in
the repository, in \file{scripts/DuplicateProbe.lean} and
\file{scripts/PlantedControls.lean}; they are outside every library target, so they cost
nothing at build time. Rerunning the probe after a refactoring is cheap and it is the only
mechanical defence against the drift that the \TauCeti{} roadmap warns about: a standard
notion said in the project's own dialect grows a redundant theory of lemmas the library
already proves. The probe measures the narrow, syntactic form of that drift. The conceptual
form---reimplementing a notion with a different statement---is what the table in
Section~\ref{sec:replace} catalogues, and no automated tool replaces reading the files.
